\section{Kết luận và Hướng phát triển}

\subsection{Tổng kết Đánh giá Dự án}
Dựa trên các phân tích và thiết kế chi tiết ở các chương trước, nhóm thực hiện xin tổng kết lại tính khả thi của giải pháp mạng cho Cơ sở Bảo Lộc như sau:

\begin{itemize}
    \item \textbf{Về Kỹ thuật:} Hệ thống tuân thủ nghiêm ngặt mô hình phân cấp 3 lớp (Core-Distribution-Access), đảm bảo tính dự phòng cao (High Availability) nhờ cơ chế VRRP và Link Aggregation. Kiến trúc này giải quyết tốt bài toán về độ ổn định cho kết nối từ xa.
    \item \textbf{Về An ninh:} Việc áp dụng tư duy "Zero Trust" ngay từ khâu thiết kế (phân chia VLAN, Port Security, Firewall Fortinet) giúp giảm thiểu rủi ro tấn công nội bộ và lan truyền mã độc.
    \item \textbf{Về Tài chính:} Với tổng dự toán \textbf{896.500.000 VNĐ}, dự án đã đạt được mục tiêu tối thượng là không vượt quá ngân sách 1 tỷ đồng, trong khi vẫn giữ lại được các thiết bị lõi (Core/Firewall) chất lượng cao.
\end{itemize}

\subsection{Hạn chế và Giải pháp}
Do giới hạn về kinh phí, thiết kế hiện tại vẫn tồn tại một số điểm hạn chế chấp nhận được:
\begin{itemize}
    \item Hệ thống Wifi chưa có Controller phần cứng chuyên dụng (đang dùng Cloud Controller miễn phí của Aruba).
    \item Chưa có hệ thống IPS (Chống xâm nhập) chuyên sâu do license Fortinet khá đắt đỏ.
\end{itemize}

\subsection{Hướng phát triển trong tương lai (Giai đoạn 2)}
Khi nhà trường có thêm ngân sách hoặc nhu cầu mở rộng (Scalability), hệ thống có thể nâng cấp dễ dàng theo lộ trình:
\begin{itemize}
    \item \textbf{SD-WAN:} Triển khai SD-WAN để tối ưu hóa đường truyền VPN kết nối về Trụ sở chính, giảm chi phí thuê kênh Leased Line.
    \item \textbf{Network Monitoring (NMS):} Trang bị thêm phần mềm giám sát (như PRTG hoặc SolarWinds) để cảnh báo sự cố chủ động qua Email/SMS.
    \item \textbf{IoT Ready:} Mở rộng thêm các Switch PoE tại các khu vực ngoài trời để lắp đặt thêm Camera AI và Cảm biến môi trường.
\end{itemize}